\section*{Chapter \ref{ch:nw:networking-for-trading} --- Networking for Trading}

\begin{solution}{exo:nw:networking-for-trading:1}
$8 \times (256 + 20) / 25 = \qty{88.3}{\nano\second}$. A 64-byte frame occupies 84 bytes of line time, 672 bits:
$10^{10}/672 \approx 14.88$ million frames a second.
\end{solution}

\begin{solution}{exo:nw:networking-for-trading:2}
The low 23 bits of 233.54.12.111 are 54.12.111 (54 is below 128, so nothing is dropped), which gives the address
\texttt{01:00:5e:36:0c:6f}. Any group
with the same low 23 bits shares it, for example 224.54.12.111 or 233.182.12.111 ($182 = 54 + 128$).
\end{solution}

\begin{solution}{exo:nw:networking-for-trading:3}
$2 \times 1.1 \times 50.1 = \qty{110.2}{\giga\bit\per\second}$ at the 10-millisecond peak; $110.2/25 = 4.4$, so five
\qty{25}{\giga\bit\per\second} links (or two at 100).
\end{solution}

\begin{solution}{exo:nw:networking-for-trading:4}
With \omterm{def:nw:networking-for-trading:nagle}{Nagle's algorithm} on, the cancel is held while the new order is unacknowledged; the venue acknowledges with its next
message to the firm (the order's acceptance) or when its delayed-acknowledgement timer fires, which TCP allows to take up to
half a second. The cancel is delayed, never reordered: TCP delivers bytes in the order written. \texttt{TCP\_NODELAY} removes
the wait.
\end{solution}

\begin{solution}{exo:nw:networking-for-trading:5}
The excess is $4 \times 6 - 10 = \qty{14}{\giga\bit\per\second}$ for \qty{50}{\micro\second}: 700\,000 bits, 87\,500 bytes. The
last byte waits $87\,500 \times 8 / 10 = 70\,000$~ns, \qty{70}{\micro\second}.
\end{solution}

\begin{solution}{exo:nw:networking-for-trading:6}
January 2024: $80 / 42.4 = 1.9$; July 2026: $350 / 63.9 = 5.5$. The second grew by half, the busiest millisecond more than
fourfold: bursts became sharper. A link sized in 2024 for that year's busiest millisecond (80 million messages a second)
carries less than a quarter of the busiest millisecond of July 2026 (350 million).
\end{solution}

\begin{solution}{exo:nw:networking-for-trading:7}
About 0.32~MB with 10\% of packets in common clusters, 1.10~MB with 30\% and 1.98~MB with 50\% (means over twenty seeds). The
buffer grows roughly in proportion to the share: the common clusters' size, and so the backlog each event leaves, is
proportional to it, while independent bursts alone need 26~KB.
\end{solution}

\begin{solution}{exo:nw:networking-for-trading:8}
A one-minute peak averages over sixty thousand milliseconds; the loss happens in the few milliseconds, or less, in which the
port is offered more than \qty{10}{\giga\bit\per\second}, which the counter cannot show. Check the switch's egress discard
counters for that port and measure the traffic at millisecond and microsecond resolution (chapter~5's captures) before blaming
the server.
\end{solution}

\begin{solution}{pb:nw:networking-for-trading:1}
\begin{enumerate}
\item $8 \times 0.6 = \qty{4.8}{\giga\bit\per\second}$: 48\% of the port.
\item Each frame is 220 bytes of line time, 1\,760 bits: $4.8 \times 10^9 / 1\,760 \approx 2.73$ million frames a second.
\item $8 \times 2.5 = \qty{20}{\giga\bit\per\second}$, twice the \omterm{def:nw:networking-for-trading:ser}{line rate}: a utilisation of 2.
\item One millisecond of excess in a minute moves the one-minute average by less than 0.002\%.
\item $(20 - 10) \times 10^9 \times 10^{-3} / 8 = 1.25$~MB.
\item $1\,250\,000 / 220 \approx 5\,682$ frames.
\item $8 \times 1.25 \times 10^6 / 10^{10} = \qty{1}{\milli\second}$.
\item It drains at $10 - 4.8 = \qty{5.2}{\giga\bit\per\second}$: $10^7 / 5.2 \times 10^9 \approx \qty{1.92}{\milli\second}$.
\item 1.25~MB.
\item $6\,000\,000 / 16 = 375\,000$ bytes.
\item $1\,250\,000 - 375\,000 = 875\,000$ bytes, about 3\,977 frames.
\item The burst offers $20 \times 10^9 \times 10^{-3} / 1\,760 \approx 11\,364$ frames: 35\% are dropped.
\item Yes: 6~MB is more than 1.25~MB, if the port may borrow the whole block.
\item \textbf{Named result.} Zero loss needs 1.25~MB of buffer for this port; with its sixteenth of the shared block, 375~KB, the
event drops about 3\,977 frames, 35\% of the burst.
\item $6 \times 10^6 \times 8 / 10^{10} = \qty{4.8}{\milli\second}$ of the same burst.
\item No: 20 is below 25.
\item No: the copies arrive at the same time, on ports of the same block, so each gets the same share and both lose the same
kind of frames at the same moment; the lines should go through different switches or blocks.
\item The switch's egress discard counter for the port (per queue); the handler's gap counter, with gaps filled by neither line.
\item Gaps on both lines go to the retransmission server if it has them within its window and rate limits, otherwise to the next
snapshot; either way the book is stale meanwhile.
\item Because the market event that makes one feed burst makes every feed burst, and a shared buffer is empty only when the
bursts are not.
\end{enumerate}
\end{solution}

\begin{solution}{iq:nw:networking-for-trading:1}
The time to clock a frame onto the wire: $8(\ell_{\mathrm f}+20)/R_{\mathrm L}$, 67.2~ns for the smallest frame at 10, 26.9 at
25. It is paid at the sender and again at every store-and-forward hop, independently of distance, so it is a large part of a
short path.
\iqlookfor{the formula with the 20 bytes of overhead, and that it recurs per hop.}
\end{solution}

\begin{solution}{iq:nw:networking-for-trading:2}
Market data goes one to many and must not slow down for anyone: multicast sends one copy that the network replicates, and UDP
has no flow control; loss is handled by redundant lines, retransmission and snapshots. Orders go one to one and must arrive
exactly once and in order, which TCP guarantees; the session layer builds on it.
\iqlookfor{one-to-many against one-to-one, and where loss recovery lives in each case.}
\end{solution}

\begin{solution}{iq:nw:networking-for-trading:3}
Nagle holds new small data while earlier data is unacknowledged; the receiver may delay its acknowledgement up to a timer. Together
a second small order can wait for the timer. Set \texttt{TCP\_NODELAY}, write each message in one call, keep the connection warm,
and acknowledge promptly where you control the receiving side.
\iqlookfor{the interaction of the two mechanisms, not only the option's name.}
\end{solution}

\begin{solution}{iq:nw:networking-for-trading:4}
Upstream of the server: the switch port toward it (egress discards during \omterm{def:nw:networking-for-trading:burst}{microbursts}), the layer-1 or switch path of both lines
if they share it, the venue handoff. On the server: the kernel's receive-buffer drops, which the card does not count. Look at
per-port discard counters and at a capture timestamped at the tap.
\iqlookfor{a common cause for both lines and a search that follows the packet path.}
\end{solution}

\begin{solution}{iq:nw:networking-for-trading:5}
It reads IGMP membership reports and sends each group only to the ports that joined it. Without it every multicast frame is
flooded to every port: links and servers carry and discard traffic they did not ask for, and \omterm{def:nw:networking-for-trading:burst}{microbursts} land on ports that
should have been quiet.
\iqlookfor{flooding as the failure, and its cost in bandwidth and CPU.}
\end{solution}

\begin{solution}{iq:nw:networking-for-trading:6}
From the correlated bursts, not the average: measure the feeds' joint arrival at microsecond resolution on busy days, compute the
backlog of the fluid bound or replay the capture through a queue model, take a high quantile of the peak backlog, and compare it
with the buffer the port can actually get when its neighbours are busy too. If it does not fit, raise the port's speed or spread
the feeds.
\iqlookfor{correlation across feeds, shared buffers, and a measured rather than assumed burst.}
\end{solution}
